Spatial covariance and within-sensor temporal order describe different aspects
of IoT behavior. We compare local correlation-spectrum, permutation and sample
entropy with conventional spatial and temporal statistics under joint block-bootstrap
and graph-conditioned diffusion references. A frozen extension evaluates
360 fault realizations on 20 simulation/recording blocks with three model seeds,
using one synthetic benchmark and two semi-synthetic real-data benchmarks.
Aligned diagnostics verify that time reordering can preserve covariance while
changing temporal entropy; autocorrelation also detects this distinction.
Results depend on support, reference and calibration. On Intel, bootstrap temporal
entropy achieves recall 0.389 versus 0.000 for the conventional temporal
scan, with respective background exceedance 0.103 and 0.000.
Adding entropy to the conventional combined scan increases synthetic bootstrap
recall by 0.003; other dataset/reference combinations show no recall increase.
Shared-support fidelity exposes substantial diffusion entropy miscoverage.
No consistent overall advantage is established for either entropy or the
generative reference. These exploratory results distinguish feature information
from estimator availability and decision performance, without claiming equivalence.
